% ============================================================
% f2014_guangdong_science_q15_parallelogram.tex
% 2014 广东卷（理科）第 15 题：平行四边形 ABCD 中，E 在 AB 上，
% EB=2AE，DE 与 AC 交于点 F，求 △CDF 与 △AEF 的面积比。
% 来源 PDF：普通高考/2014/2014广东理.pdf
% 原页：第 1 页（A3 横排）第 2 栏，印面 877，第 15 题题干下方
% 说明：AE:EB=1:2，故 AE=AB/3；图中 A、E、B 依次在底边上，
% D 在 A 上方偏右，C 在 B 上方偏右，DE 近于竖直，AC 为对角线。
% ============================================================
\begin{tikzpicture}[>=Stealth, line width=0.7pt, font=\small,
  x=1cm, y=1cm, line cap=round, line join=round]

  \coordinate (A) at (0,0);
  \coordinate (E) at (1.3333,0);
  \coordinate (B) at (4,0);
  \coordinate (D) at (0.8,3);
  \coordinate (C) at (4.8,3);
  \coordinate (F) at (1.2,0.75);

  \draw (A) -- (B) -- (C) -- (D) -- cycle;
  \draw (D) -- (E);
  \draw (A) -- (C);

  \fill (A) circle (1.1pt) node[below left] {$A$};
  \fill (E) circle (1.1pt) node[below] {$E$};
  \fill (B) circle (1.1pt) node[below right] {$B$};
  \fill (D) circle (1.1pt) node[above left] {$D$};
  \fill (C) circle (1.1pt) node[above right] {$C$};
  \fill (F) circle (1.1pt) node[right] {$F$};
\end{tikzpicture}
